% CODE-FREE literal layout draft. Never compiled or rendered in this task.
% All coefficient data and every residual family are printed below.
\documentclass[10pt,letterpaper]{article}
\usepackage[T1]{fontenc}
\usepackage{lmodern,amsmath,amssymb,multicol,geometry}
\geometry{margin=0.43in}
\pagestyle{empty}
\setlength{\parindent}{0pt}
\setlength{\parskip}{2pt}
\setlength{\columnsep}{13pt}
\setlength{\abovedisplayskip}{2.5pt}
\setlength{\belowdisplayskip}{2.5pt}
\setlength{\abovedisplayshortskip}{1.5pt}
\setlength{\belowdisplayshortskip}{1.5pt}
\setlength{\multicolsep}{3pt}
\newcommand{\NN}{\mathbb N}
\newcommand{\PP}{\mathcal P}
\newcommand{\QQ}{\mathcal Q}
\newcommand{\HH}{\mathcal H}
\begin{document}
\fontsize{9}{10.2}\selectfont
\begin{center}
{\fontsize{11}{12}\selectfont\bfseries An indexed C222 Diophantine equation}\par
\textit{Complete polynomial specification; unreviewed text-only layout draft}
\end{center}
\vspace{-4pt}
\begin{multicols}{2}
\textbf{Free coordinates.}
Every free coordinate is in $\NN=\{0,1,2,\ldots\}$:
\[
\begin{gathered}
h,e,g,\beta,F,x,N,\ell,z_0,z_t,z_B,\quad(\xi_i)_{i=0}^8,\\
(D_i)_{i=0}^3,\quad(P_j)_{j=0}^{10},\quad
(r_j,s_j,w_j)_{j=0}^{10},\quad(\zeta_j)_{j=0}^{17}.
\end{gathered}
\]
Each $\zeta_j$ is the 27-tuple specified below. All other names are
polynomial abbreviations; no further free coordinates occur. Put
\[
\begin{aligned}
t&=16+h,&y&=F+1+e,&q&=2+g,&B&=1+\beta,\\
(C,a,b,c,d,S,T,u,v)&=(1+\xi_i)_{i=0}^8,\\
H&=2+C+a+b+c+d+S+T+u+v.
\end{aligned}
\]

\textbf{The numerical calculus.}
For every use of these polynomial functions, put
\[
\begin{aligned}
I(a,b)&=2(a+b)(a+b+1)+4b,\\
M(a,b)&=I(a,b)+1,&A(a,b)&=I(a,b)+2,\\
E(a,b)&=I(a,b)+3,&\nu(a)&=I(a,1),\\
X(a,b)&=\nu(A(a,\nu(b))),&
\Theta_a(z)&=\nu(A(a,z)).
\end{aligned}
\]
The first eight code polynomials are
\[
\begin{aligned}
c_1&=X(a,E(a,b)),\\
c_2&=I(A(a,I(b,c)),I(A(a,b),A(a,c))),\\
c_3&=I(E(a,b),I(M(c,a),M(c,b))),\\
c_4&=I(A(a,A(b,c)),A(b,A(a,c))),\\
c_5&=I(E(c,a),I(M(c,b),M(a,b))),\\
c_6&=I(X(a,A(a,b)),b),\\
c_7&=I(E(c,a),I(E(c,b),A(M(a,b),E(a,b)))),\\
c_8&=I(I(I(c,d),a),I(I(a,c),I(b,c))).
\end{aligned}
\]
The complete ninth code is defined by
\[
\begin{aligned}
R(b,k,z)&=A(b,I(z,\Theta_b(I(z,\nu(M(b,k)))))),\\
Q(a,b,k)&=A(k,I(M(k,a),M(k,b))),\\
D(a,b,k,z)&=I(Q(a,b,k),R(b,k,z)),\\
C_*(a,b,k,z)&=A\bigl(b,I(M(b,k),\\[-2pt]
&\hspace{39pt}\Theta_a(I(z,\Theta_a(D(a,b,k,z)))))\bigr),\\
V(a,b,c)&=I(M(a,b),A(c,M(a,b))),\\
K(a,b,c)&=A(c,I(C_*(a,b,c,c),\nu(R(b,c,c)))),\\
c_9&=\nu(I(V(a,b,c),K(a,b,c))).
\end{aligned}
\]
Here each $c_i$ is evaluated at $(a,b,c,d)$ and ignores its unused
arguments. Define the following polynomials:
\[
\begin{aligned}
W_z&=1+Fz,\\
J&=\sum_{i=1}^9c_i+a+b+I(a,b)+A(a,b)+1,\\
G&=((c-a)^2(c-b)^2-(d+1)^2)^2,\\
\HH&=(T-S)^2\prod_{i=1}^8(TW_{c_i}-S)^2\\
&\quad\cdot\bigl((TW_{c_9}-S)^2+G\bigr)\\
&\quad\cdot(TW_b-SW_aW_{I(a,b)})^2\\
&\quad\cdot(TW_{A(a,b)}-SW_b)^2,\\
L_i&=\HH+(x-S-(1+(i+1)y)u)^2\\
&\hspace{32pt}+(x-T-(1+(i+2)y)v)^2.
\end{aligned}
\]
Only $L_C$ is used below. On a vanishing guarded branch, $G=0$
certifies $c\ne a,c\ne b$ with its natural witness $d$.

\textbf{One fully printed power certificate.}
In this paragraph every letter of $\zeta$ is local to this copy,
independently of all front coordinates and other copies:
\[
\begin{gathered}
\zeta=(W,A,\tau,Z,X,Y,U,V,S,R,B,a_0,y_0,b_0,f,d,\\
b_+,b_-,s_+,s_-,r_+,r_-,x_+,x_-,g,w_n,w_k).
\end{gathered}
\]
Let $\mathcal N(a,x,y)=x^2+y^2-1-a^2y^2$. For arguments $(n,k,m)$
let $(R_1,\ldots,R_{17})$ be the following ordered tuple:
\[
\begin{gathered}
\bigl(A-a_0-2,\ Y-k-y_0,\ B-b_0-2,\\
\mathcal N(A,X,Y),\ \mathcal N(A,U,V),\ \mathcal N(B,S,R),\\
B-1-4Yf,\ V-Y^2(d+1),\\
B+Ub_+-A-Ub_-,\ S+Us_+-X-Us_-,\\
R+4Yr_+-k-4Yr_-,\\
X+Yn+\tau x_+-YA-m-\tau x_-,\\
2An-\tau-n^2-1,\ \tau-m-g-1,\\
W-n-w_n,\ W-k-w_k,\ \mathcal N(W+1,A,WZ)\bigr).
\end{gathered}
\]
Set $\PP(n,k,m;\zeta)=\sum_{i=1}^{17}R_i^2$. Every slack and
each $+$ or $-$ quotient coordinate in this tuple is natural.

\textbf{All eighteen argument slots.}
The following vectors and ordered concatenation give every slot:
\[
\begin{aligned}
(k_i)_{i=0}^3&=(t,t+1,y+1,x),\\
(U_j)_{j=0}^{10}&=(F^2+t,C,a,b,c,d,S,T,u,v,J),\\
(\sigma_j)_{j=0}^{10}&=(0,0,1,1,1,1,2,2,3,3,1),\\
\mathcal T&=\bigl((2(t+1),t,F),(B,x,N),(BH,x,q),\\
&\qquad(q,k_i,D_i)_{i=0}^3,\ (q+1,U_j,P_j)_{j=0}^{10}\bigr).
\end{aligned}
\]
Write $(n_j,k'_j,m_j)=\mathcal T_j$ for $0\le j\le17$ in this
printed order, and pair it with $\zeta_j$. These are polynomial
substitutions, with no variable exponent.

\textbf{The one equation.}
\[
\boxed{\begin{aligned}
\QQ={}&\sum_{j=0}^{17}\PP(n_j,k'_j,m_j;\zeta_j)\\
&+(x-F-1-(y+1)z_0)^2\\
&+(x-1-(1+(t+1)y)z_t)^2\\
&+(Bz_B-C-F^2-1)^2+(L_C-N\ell)^2\\
&+(P_0-D_0(qw_0+B)-r_0)^2\\
&+\sum_{j=1}^{10}(P_j-D_{\sigma_j}Nw_j-r_j)^2\\
&+\sum_{j=0}^{10}(D_{\sigma_j}-r_j-s_j-1)^2=0.
\end{aligned}}
\]
The final sum imposes all eleven natural remainder bounds; the preceding
rows are the carrier digit, ten range digits and four outer residuals.
Expanding only fixed compositions, sums and products defines an ordinary
integer-coefficient polynomial. The entire coefficient recipe is printed here.

The hand census is 554 natural coordinates and 332 squared residuals.
For positive unknowns, substitute $z=p_z-1$, $p_z\ge1$, independently
for every free coordinate $z$. No coordinate is shared between private
power copies. The complete ordinary correctness proposal is separate;
this layout has not been compiled, rendered, visually reviewed, or accepted.
\end{multicols}
\end{document}
